\setupcolors[state=start]
\setupcolor[hex]
\definecolor[darkblue][h=003366]
\definecolor[darkgray][h=e6e6e6]
\definecolor[redline][h=d9534f]
\definecolor[blueline][h=5bc0de]
\definecolor[greenline][h=5cb85c]

\setuphead[section][style=\bfa, color=darkblue, before=\blank, after=\blank]
\setuphead[subsection][style=\bf, color=black, before=\blank, after=\blank]
\setuppapersize[A4][A4]
\setupbodyfont[11pt]
\setuplayout[width=middle, height=middle, topspace=2cm, bottomspace=2cm]
\mainlanguage[es]

\starttext
\title{\bf Comparación de OPP\_v5\_1\_mono\crlf y OPP\_v5\_2\_split en pruebas funcionales y de stress}
\blank[big]

\startsection[title={Resumen}]
\startitemize
\item {\bf Mismo problema, mismo óptimo}: en los %%NUM:feasible%% pares que los dos resuelven el
objetivo es el mismo, y los %%NUM:infeasible%% infactibles lo son en los dos.
\item {\bf En instancias chicas da lo mismo}: el monolítico es algo más rápido por décimas de
segundo, porque compila un modelo en vez de dos.
\item {\bf A escala realista solo el split prueba el óptimo}: con la misma instancia, el monolítico
deja de cerrar en 10 minutos lo que el split resuelve en segundos, y con 40 minutos tampoco lo
logra.
\stopitemize
\stopsection

\startsection[title={Objetivo y alcance}]
Este documento compara las ejecuciones de los tests funcionales y de stress de OPP\_v5\_1\_mono
(formulación monolítica) y OPP\_v5\_2\_split (formulación en dos etapas) sobre los mismos
escenarios. Los dos modelos definen el mismo problema de optimización (ruteo, asignación y
secuenciación, con velocidades heterogéneas de pickers y cadena de frío) y tienen el mismo valor
óptimo.

Los datos son los CSV consolidados de cada modelo:
\startitemize
\item \type{results/opp_split/functional/} — %%NUM:split_tests%% tests, con la Etapa 1 en
secuencia.
\item \type{results/opp_mono/functional/} — %%NUM:mono_tests%% tests.
\item \type{results/comparability/compare_split_vs_mono.csv} — la comparación de las
soluciones guardadas de los dos modelos, test por test.
\stopitemize

La suite de stress (sección \in[sec:stress]) usa \type{results/opp_split/stress/} y
\type{results/opp_mono/stress/}: los mismos %%NUM:stress_total%% escenarios en los dos
modelos, con un tope de 40 minutos por test.
\stopsection

\startsection[title={Qué tests se comparan}]
Se comparan solo los tests {\em comparables}: los que OPP\_v5\_2\_split, OPP\_hc\_4 y OPP\_aco\_4
resuelven sobre la misma instancia y con el mismo resultado esperado, según el registro
\type{results/comparability/comparability_split_vs_heuristics.csv} (reglas en
\type{docs/specs/comparabilidad_split_heuristicas.md}). Es la misma base que la comparación de
OPP\_v5\_2\_split con las heurísticas, así que los dos documentos hablan de los mismos tests.
OPP\_v5\_1\_mono corre los mismos tests que el split.

Dentro de esa base, un par es comparable si los dos modelos lo resuelven (se comparan objetivo,
solución y tiempos) o si los dos lo declaran infactible (se compara el estado).

%%TABLE:coverage%%

Quedan fuera %%NUM:excluded%% tests. El 37a sería comparable entre mono y split (los dos
modelan la espera voluntaria), pero se deja fuera para mantener la misma base que la comparación
con las heurísticas:

%%TABLE:excluded%%
\stopsection

\startsection[title={Equivalencia del valor objetivo}]
De los %%NUM:feasible%% pares resueltos, {\bf %%NUM:same_obj%% tienen el mismo valor objetivo}. La
diferencia máxima es $%%NUM:max_diff%%$ (casos %%NUM:diff_cases%%), dentro de la tolerancia de
SCIP. Los %%NUM:infeasible%% pares infactibles lo son en los dos modelos.

Esto confirma empíricamente que las dos formulaciones tienen el mismo valor óptimo, también en
los tests de cadena de frío, velocidad y ruta más corta.
\stopsection

\startsection[title={Soluciones: mismo óptimo, no siempre la misma solución}]
Los dos modelos alcanzan el mismo objetivo, pero no siempre con la misma solución. En
%%NUM:same_solution%% de los %%NUM:feasible%% pares devuelven la misma; en los otros
%%NUM:alternative%% el monolítico devuelve un {\em óptimo alternativo}: otras rutas
(%%NUM:alt_routes%% pares), otra asignación de lotes a pickers o de órdenes a lotes
(%%NUM:alt_assign%%), o las dos cosas.

Es lo esperable. El objetivo no tiene un término de tiempo de ruta, así que el monolítico puede
elegir cualquier ruta que no empeore las penalizaciones; en el split, la Etapa 1 siempre produce la
ruta mínima, porque su objetivo es minimizar el tiempo de ruta. Las asignaciones distintas son
empates del objetivo que cada solver rompe a su manera. La columna {\em Solución de mono} de la
sección \in[sec:detalle] dice en qué difiere cada par.
\stopsection

\startsection[title={Tiempos de los pares resueltos}]
Sobre los %%NUM:feasible%% pares resueltos. \type{total_time} es la ejecución completa del test,
incluida la compilación de ZIMPL; \type{model_time} es solo el solver (en el split, la Etapa 2).

%%TABLE:times%%

En total, el monolítico tardó %%NUM:sum_ratio%% veces más que el split, pero en la mediana la
diferencia se invierte: %%NUM:median_mono%%~s contra %%NUM:median_split%%~s. En %%NUM:mono_faster%%
de los %%NUM:feasible%% pares el monolítico fue {\bf más rápido} (entre %%NUM:mono_faster_min%%× y
%%NUM:mono_faster_max%%×): son instancias chicas, donde compilar y resolver un solo modelo cuesta
menos que compilar dos. La suma la dominan las instancias de grilla, donde el monolítico llega a
tardar %%NUM:top_ratio%% veces más (caso %%NUM:top_case%%).

En \type{model_time} la diferencia es mayor (%%NUM:model_sum_ratio%% veces en la suma,
%%NUM:model_median_ratio%% en la mediana), porque el monolítico resuelve un MIP con más variables y
restricciones que cualquiera de las dos etapas.
\stopsection

\startsection[title={Visualización comparativa}]

\startsubsection[title={Tiempo total: monolítico contra dos etapas}]
Cada punto es un par resuelto. Los puntos por encima de la diagonal son instancias donde el
monolítico fue más lento.

\placefigure[force][fig:scatter]{Tiempo total de cada par resuelto: monolítico contra split (escala logarítmica en los dos ejes).}
{%%FIG:scatter%%}
\stopsubsection

\startsubsection[title={Cociente de tiempos por par}]
Barras ordenadas por el cociente mono / split. Las que quedan por debajo de 1 son instancias
donde el monolítico fue más rápido. Se identifican los casos en que el monolítico tarda más de
10 veces.

\placefigure[force][fig:ratio]{Tiempo total del monolítico dividido el del split, por par resuelto (escala logarítmica).}
{%%FIG:ratio%%}
\stopsubsection

\startsubsection[title={Tiempo del solver según el tamaño}]
Tiempo del solver contra el tamaño de la instancia, medido como lotes × ubicaciones, para los
%%NUM:sized%% pares cuyo CSV registra ese tamaño. La separación entre las dos nubes crece con el
tamaño: la descomposición escala mejor.

\placefigure[force][fig:size]{Tiempo del solver contra el tamaño de la instancia (escala logarítmica en los dos ejes).}
{%%FIG:size%%}
\stopsubsection
\stopsection

\startsection[title={Detalle de los pares resueltos}, reference=sec:detalle]
%%TABLE:detail%%
\stopsection

\startsection[title={Tests infactibles}]
Los %%NUM:infeasible%% pares que los dos modelos declaran infactibles.

%%TABLE:infeasible%%
\stopsection

\startsection[title={Escala: la misma instancia realista en grillas crecientes}, reference=sec:escala]
Los tests funcionales son instancias chicas: el solver tarda milisegundos y el tiempo total lo
pone la compilación de ZIMPL, así que ahí los dos modelos son intercambiables. Para ver qué pasa
cuando la instancia crece, el estudio 08 resuelve con los dos modelos {\em la misma} instancia del
perfil realista (26 órdenes por lote, pesos 4/5, dos bandas de velocidad, dos lotes por picker)
en grillas de 3×3 a 7×7 y con 1 a 6 lotes, con un tope de %%NUM:scale_cap_min%% minutos. En 4×4
la instancia es infactible (no entran las ubicaciones) y no se grafica.

\placefigure[force][fig:escala]{Tiempo hasta probar el óptimo, en la misma instancia realista, según la grilla y la cantidad de lotes (escala logarítmica, la misma en los dos paneles). Estudio 08.}
{\startcombination[2*1]
{%%FIG:scale_split%%}{OPP\_v5\_2\_split}
{%%FIG:scale_mono%%}{OPP\_v5\_1\_mono}
\stopcombination}

El contraste es de forma, no de grado. En el split las curvas crecen despacio: con hasta 4 lotes
resuelve cualquier grilla en menos de 10~s. En el monolítico se disparan: con 3 lotes pasa de 9~s
en 5×5 a 139~s en 6×6 y a no terminar en 7×7, y en %%NUM:scale_capped%% no prueba el óptimo en
%%NUM:scale_cap_min%% minutos. Entre los puntos que los dos cierran, el monolítico tarda hasta
%%NUM:scale_max_ratio%% veces más (%%NUM:scale_max_case%%); donde no cierra, la relación es de
al menos %%NUM:scale_bound_ratio%% veces (%%NUM:scale_bound_case%%, que el split resuelve en
%%NUM:scale_bound_split%%~s).

La razón es estructural. El estudio registra el tamaño del MIP que resuelve SCIP: en el
monolítico, el modelo entero; en el split, la Etapa 2, porque la Etapa 1 son MIPs chicos de ruteo,
uno por lote, que se resuelven en pocos segundos.

%%TABLE:model_size%%

El MIP del monolítico crece con la grilla, porque incluye el ruteo de todos los lotes; el de la
Etapa 2 no depende de la grilla, porque recibe las rutas ya resueltas y solo decide asignación y
secuencia. Por eso el split escala con la grilla y el monolítico no.

¿Es cuestión de darle más tiempo? El estudio 11 repite esos puntos, y la instancia exacta del
test de stress 04, con 40 minutos de tope:

%%TABLE:long_cap%%

El monolítico {\em encuentra} soluciones muy buenas (en el stress 04, a %%NUM:t04_best%%~\% del
óptimo), pero no puede {\em probar} que lo son: su cota dual casi no se mueve, y en el stress 04
reporta un gap de %%NUM:t04_gap%%~\% a los 40 minutos. El split prueba el óptimo de esa instancia en
%%NUM:t04_split%%~s. Sin el óptimo del split, la solución del monolítico no se podría distinguir de
una mala.

Los estudios 08 y 11 corrieron con la máquina compartida con otras corridas, así que sus tiempos
absolutos están inflados; lo que se compara es la relación entre los dos modelos sobre la misma
instancia, que corrieron uno después del otro. El detalle está en el documento de exploración del
monolítico.
\stopsection

\startsection[title={Suite de stress},reference=sec:stress]
Los tests de stress llevan los dos modelos a instancias de tamaño real: el perfil realista (26
órdenes por lote, dos bandas de velocidad, pesos 4/5) en grillas de 7×7 a 16×16 y cuadrillas de
hasta 24 lotes, más tres escenarios heredados de OPP\_v4. El split corre la Etapa 1 con seis
procesos. Su tiempo es la duración completa del test y es un techo: si la notebook se suspende
durante la corrida, el tiempo suspendido cuenta en la duración pero no en el tope.

Como el monolítico no termina ninguno, no hay pares resueltos en los dos: la comparación de stress
es solo de estado. En el 01 (100×100) el monolítico figura {\em sin resultado}: el proceso del test
terminó sin devolver un estado, ni solución ni timeout.

%%TABLE:stress%%

El monolítico termina %%NUM:stress_mono_ok%% de los %%NUM:stress_total%% tests; el split,
%%NUM:stress_split_ok%%, ninguno en más de %%NUM:stress_split_max%% minutos. En los tests
funcionales las dos formulaciones son intercambiables salvo en las grillas; a este tamaño, solo
el split da una respuesta. Los tests que tampoco termina el split son los del perfil realista
con 8 lotes o más (12 a 14 y 16 a 19), la grilla de 16×16 (08) y dos de los heredados (01 y
03).
\stopsection

\startsection[title={Conclusiones}]
{\bf Para instancias de tamaño real, el split es claramente superior: es la única de las dos
formulaciones exactas que prueba el óptimo.} El monolítico solo es competitivo en instancias
chicas, donde la diferencia es de décimas de segundo.

\startitemize[n]
\item {\bf Mismo problema, mismo óptimo.} Los %%NUM:feasible%% pares resueltos dan el mismo
objetivo (diferencia máxima $%%NUM:max_diff%%$, redondeo del solver) y los %%NUM:infeasible%%
infactibles lo son en los dos. En %%NUM:alternative%% pares el monolítico llega a ese óptimo con
otra solución (otras rutas u otra asignación), porque el objetivo no penaliza la longitud de la ruta.
\item {\bf Instancias chicas: da lo mismo.} En los tests funcionales el solver tarda
milisegundos; el monolítico es algo más rápido en %%NUM:mono_faster%% de %%NUM:feasible%% pares
(mediana %%NUM:median_mono%%~s contra %%NUM:median_split%%~s) solo porque compila un modelo en
vez de dos.
\item {\bf En cuanto la instancia crece, el monolítico se queda atrás.} Ya en las grillas de los
tests funcionales tarda hasta %%NUM:top_ratio%% veces más. Con el perfil realista (sección
\in[sec:escala]) su tiempo se dispara con la grilla: con 3 lotes en 7×7 o 4 lotes en 6×6 no cierra
en 10 minutos lo que el split resuelve en menos de 10~s.
\item {\bf Más tiempo no lo arregla.} Con 40 minutos, el monolítico encuentra soluciones a menos
de 1~\% del óptimo pero no lo prueba: en el stress 04 reporta un gap de %%NUM:t04_gap%%~\% sobre
una instancia que el split cierra en %%NUM:t04_split%%~s.
\item {\bf A escala de stress, solo responde el split.} En los %%NUM:stress_total%% tests de stress
el monolítico no termina ninguno en 40 minutos; el split termina %%NUM:stress_split_ok%%, hasta 6
lotes en 15×15. Con 8 lotes o más, o 16×16, tampoco el split termina: ahí la alternativa son las
heurísticas, que se comparan con el split en un documento aparte.
\stopitemize
\stopsection

\stoptext
